\section{Data Structures and Algorithms}\label{sec:dsa}

This section describes how a task is represented, in the database and in memory, and the algorithms behind the
program's operations.

\subsection{The task}

A task is an instance of the JPA entity \texttt{Task} in package \texttt{entities}, mapped onto the table
\texttt{tasks}. Table~\ref{tab:task} lists its fields. The two length limits are public constants of the class,
\texttt{Task.MAX\_TITLE\_CHARACTERS} (100) and \texttt{Task.MAX\_DESCRIPTION\_CHARACTERS} (1000), so that the
service and the entity's column definitions use the same numbers. The fields also carry Bean Validation
annotations (\texttt{@NotBlank}, \texttt{@NotNull}, \texttt{@Size}); they are not where the rules are enforced, which
is the service, but a last check that stops a task that slipped past the service from reaching the database.

\begin{table}[!t]
  \caption{The fields of a task}\label{tab:task}
  \centering\footnotesize
  \begin{tabularx}{\linewidth}{@{}L{0.2\linewidth}L{0.16\linewidth}L{0.25\linewidth}Y@{}}
    \toprule
    Field & Java type & Column & Rule \\
    \midrule
    \texttt{id} & \texttt{UUID} & \texttt{id char(36)}, primary key & Given by Hibernate on the first save; never
      changed, so it has no setter. \\
    \texttt{title} & \texttt{String} & \texttt{title varchar(100)} & Not blank; trimmed; at most 100 characters. \\
    \texttt{description} & \texttt{String} & \texttt{description varchar(1000)} & Trimmed; at most 1000
      characters; empty, never null, when there is none. \\
    \texttt{status} & \texttt{TaskStatus} & \texttt{status varchar(16)} & A new task is Pending. \\
    \texttt{createdAt} & \texttt{Instant} & \texttt{created\_at timestamp} & Set when the task is added; never
      written again. \\
    \texttt{updatedAt} & \texttt{Instant} & \texttt{updated\_at timestamp} & Equal to \texttt{createdAt} at first;
      moved on by every update and status change. \\
    \bottomrule
  \end{tabularx}
\end{table}

Every column is \texttt{not null}. The identifier is a random UUID rather than a counter, stored as 36 characters,
and the two times are \texttt{Instant}s, moments in UTC that the screen shows in the computer's own time zone.

\subsection{The status}

The status is the enum \texttt{TaskStatus} with three constants, \texttt{PENDING}, \texttt{IN\_PROGRESS} and
\texttt{COMPLETED}. Each carries the words the screen shows for it (``Pending'', ``In Progress'', ``Completed''),
returned by \texttt{getLabel()} and by \texttt{toString()}, so a list or a choice box of statuses shows the words
without any conversion. The database keeps the constant's name (\texttt{@Enumerated(EnumType.STRING)}), not its
position, so reordering the constants or changing a label never changes what a saved task means.

\subsection{The table}

The table is created by the first Flyway migration, Listing~\ref{lst:v1}, which is the whole schema of the program.

\source[language=SQL, morekeywords={char,varchar,timestamp}]{src/main/resources/db/migration/V1__tasks.sql}
  {The migration \texttt{V1\_\_tasks.sql}}{lst:v1}

\subsection{The counts}

The number of tasks in each status is returned by the service as a record,
\texttt{TaskCounts(long pending, long inProgress, long completed)}, with two methods of its own: \texttt{total()}, the
number of tasks in all, and \texttt{of(TaskStatus)}, the count for one status, which lets the screen go through
\texttt{TaskStatus.values()} instead of naming each status.

\subsection{Collections in the interface}

The window shows the tasks in a JavaFX \texttt{ListView<Task>}, whose rows are drawn by a cell class of the
program's own, \texttt{TaskCell}. Its items are an observable list, which the list redraws whenever it changes. The
screen keeps no copy of the tasks to filter or count itself: after every change, and whenever the search text or the
filter changes, it asks the service again for the matching tasks and the counts, and replaces the list's items with
the answer, choosing the same task again if it is still listed. What the window shows is therefore always what is
saved, including what a second copy of the program has changed, once the list is read again. Times are shown in the
computer's own time zone, as ``Today, 21:40'', ``Yesterday, 09:05'', ``16 Nov, 14:00'' or, in another year,
``16 Nov 2025, 14:00''.

\subsection{Algorithms}

\paragraph*{Validation} Every save of a title and a description goes through one private method of the service,
\texttt{set}. It trims both texts with \texttt{String.strip()}, which removes spaces, tabs and line breaks around
them, and turns a missing description into an empty one. It then refuses, in this order, an empty title, a title over
100 characters and a description over 1000, each with its own message. The limits are counted after trimming, so
spaces typed around a text never push it over. An update checks before this that a task was chosen and still exists,
and after it that a status was given.

\paragraph*{Listing and ordering} The list is every task, newest first. The repository declares the method
\texttt{findAllByOrderByCreatedAtDesc()}, and Spring Data derives the query from its name.

\paragraph*{Search and filter} Searching and filtering are one database query (Listing~\ref{lst:search}). The service
trims the search text, lowers its case, and wraps it as a \texttt{like} pattern, \texttt{\%text\%}, matched against
the lowered title and the lowered description. The characters \texttt{\%} and \texttt{\_}, which are wildcards in a
\texttt{like} pattern, and \texttt{!}, chosen as the escape character, are escaped first, so that every character the
user types is matched as it is. A blank search becomes \texttt{\%\%}, which matches every task, and a status of
\texttt{null} means every status, so the same query serves the plain list, a search, a filter and both together.
Case is ignored only for the letters A to Z, because SQLite's \texttt{lower()} function folds ASCII letters only.

\source[firstline=28, lastline=42, gobble=4]{src/main/java/io/mcs/taskmanager/repositories/TaskRepository.java}
  {The search query in \texttt{TaskRepository}}{lst:search}

\paragraph*{Counts per status} The service counts each status with the derived query
\texttt{countByStatus(TaskStatus)}, three queries in one read-only transaction. The counts are always over every task,
whatever the search or filter shows.

\paragraph*{A task that no longer exists} Every operation on one task (opening it, updating it, changing its status
and deleting it) first looks it up through one private method, \texttt{find}. A missing identifier is refused with ``Choose a task in
the list first.'', and a task the database no longer has with ``This task no longer exists: it may have been deleted
in another copy of the Task Manager.'' Because each service method is one transaction and the
database begins every transaction by taking its write lock (\texttt{transaction\_mode=IMMEDIATE} in the connection
URL), the other copy of the program cannot delete the task between this lookup and the save that follows it; its
delete waits until the save is done. How this was found and fixed is told in Section~\ref{sec:git}.

\paragraph*{Cost} With $n$ tasks, listing, searching and counting each read the whole table, $O(n)$: a pattern that
starts with \texttt{\%} cannot use an index, and the list is short enough that none is needed. Finding one task by its
identifier uses the index SQLite keeps on the primary key, $O(\log n)$.
